%%
%% ACM Conference Paper — Two-Column (sigconf)
%% Title: Weaponized Weights: A Cryptographic Provenance and Binary
%%        Inspection Framework for Securing ML Model Artifacts in
%%        Enterprise MLOps Pipelines
%%
%% Authors: Sunil Gentyala, Rakesh Prakash
%%
%% Compile: pdflatex acm_paper && bibtex acm_paper && pdflatex acm_paper (x2)
%%

\documentclass[sigconf,review]{acmart}

%% -----------------------------------------------------------------------
%% ACM rights block — update before camera-ready submission
%% -----------------------------------------------------------------------
\copyrightyear{2026}
\acmYear{2026}
\setcopyright{acmcopyright}
\acmConference[ACM CCS '26]{ACM SIGSAC Conference on Computer and
  Communications Security}{November 2026}{Salt Lake City, UT, USA}
\acmBooktitle{ACM CCS '26: ACM SIGSAC Conference on Computer and
  Communications Security, November 2026, Salt Lake City, UT, USA}
\acmDOI{10.1145/XXXXXXX.XXXXXXX}
\acmISBN{978-1-4503-XXXX-X/26/11}

%% -----------------------------------------------------------------------
%% Packages
%% -----------------------------------------------------------------------
\usepackage{booktabs}
\usepackage{listings}
\usepackage{xcolor}
\usepackage{microtype}
\usepackage{graphicx}
\usepackage{amsmath}
\usepackage{balance}

\lstset{
  basicstyle=\ttfamily\footnotesize,
  breaklines=true,
  frame=single,
  captionpos=b,
  numbers=left,
  numberstyle=\tiny,
  tabsize=2,
}

%% -----------------------------------------------------------------------
%% CCS Concepts
%% -----------------------------------------------------------------------
\begin{CCSXML}
<ccs2012>
  <concept>
    <concept_id>10002978.10002997.10002998</concept_id>
    <concept_desc>Security and privacy~Software security engineering</concept_desc>
    <concept_significance>500</concept_significance>
  </concept>
  <concept>
    <concept_id>10002978.10002997.10003003</concept_id>
    <concept_desc>Security and privacy~Malware and its mitigation</concept_desc>
    <concept_significance>300</concept_significance>
  </concept>
  <concept>
    <concept_id>10010147.10010178</concept_id>
    <concept_desc>Computing methodologies~Machine learning</concept_desc>
    <concept_significance>200</concept_significance>
  </concept>
</ccs2012>
\end{CCSXML}

\ccsdesc[500]{Security and privacy~Software security engineering}
\ccsdesc[300]{Security and privacy~Malware and its mitigation}
\ccsdesc[200]{Computing methodologies~Machine learning}

%% -----------------------------------------------------------------------
%% Keywords
%% -----------------------------------------------------------------------
\keywords{machine learning security, supply chain attack, model provenance,
  Safetensors, GGUF, pickle deserialization, MLOps, cryptographic
  verification, CI/CD security, binary inspection}

%% -----------------------------------------------------------------------
%% Authors
%% -----------------------------------------------------------------------
\author{Sunil Gentyala}
\authornote{Corresponding author.}
\email{sunil.gentyala@ieee.org}
\orcid{0000-0000-0000-0000}
\affiliation{%
  \institution{HCLTech}
  \department{Cybersecurity and AI Security Practice}
  \city{Dallas}
  \state{TX}
  \country{USA}
}

\author{Rakesh Prakash}
\email{Rakesh.Prakash@colorado.edu}
\affiliation{%
  \institution{University of Colorado Boulder}
  \department{Department of Computer Science}
  \city{Boulder}
  \state{CO}
  \country{USA}
}

%% -----------------------------------------------------------------------
%% Title
%% -----------------------------------------------------------------------
\title{Weaponized Weights: A Cryptographic Provenance and Binary Inspection
  Framework for Securing Machine Learning Model Artifacts in Enterprise
  MLOps Pipelines}

%%
%% -----------------------------------------------------------------------
\begin{document}

\begin{abstract}
The proliferation of publicly distributed machine learning model artifacts
has introduced a structural security risk that conventional enterprise
tooling is fundamentally ill-equipped to address. Weight files serialized in
PyTorch checkpoint, Safetensors, and GGUF formats are routinely downloaded
from public model repositories and admitted into production inference
pipelines without binary-level inspection or cryptographic provenance
verification. Existing static analysis, dynamic analysis, and antivirus
infrastructure have no visibility into these binary weight formats, leaving
organizations exposed to a class of supply chain attack that operates
entirely below the application layer.

We present \textit{model-provenance-guard}, a production-ready, open-source
toolkit comprising a Python binary inspector (\texttt{verify\_weights.py})
and a five-stage GitHub Actions provenance gate. The pipeline enforces
SHA-256 hash verification against an internal trusted registry, Sigstore/Cosign
cryptographic signature checking, pickle opcode scanning via
\texttt{picklescan}, Safetensors header anomaly detection, and GGUF metadata
key auditing before any artifact is promoted within a downstream pipeline.
We develop a structured, format-level threat model grounded in documented
real-world incidents and published vulnerability disclosures, and evaluate
the toolkit against eight synthetic adversarial test cases that exercise
dtype manipulation, data offset injection, magic byte substitution, and
malformed JSON header attacks. The inspector achieved zero false negatives
across all adversarial scenarios and zero false positives on benign artifacts,
demonstrating that shift-left provenance enforcement is both technically
feasible and operationally practical at the CI/CD boundary.
\end{abstract}

\maketitle

%% -----------------------------------------------------------------------
\section{Introduction}
%% -----------------------------------------------------------------------

The security community's attention toward artificial intelligence systems
has concentrated heavily on prompt injection and model output manipulation
over the past two years. Those concerns are legitimate, but their prominence
has allowed a structurally more dangerous threat to mature with insufficient
scrutiny. Adversaries are no longer confined to crafting adversarial inputs
at runtime. They are embedding threats directly inside compiled weight
artifacts that machine learning engineers download from public repositories
as a routine operational act.

The attack surface has migrated from the application layer to the binary
layer, and the tooling most enterprises rely on for security assurance lacks
any meaningful visibility there. Static application security testing (SAST)
tools analyze source code. Dynamic application security testing (DAST)
tools probe running applications. Antivirus engines match file content against
known malware signatures. None of these approaches has any concept of what
lies inside a 4-bit quantized GGUF binary, a Safetensors archive, or a
PyTorch checkpoint serialized with Python's \texttt{pickle} protocol. When
an engineer issues a \texttt{huggingface-cli download} command, the SIEM
records a successful HTTPS transfer. The endpoint detection and response
(EDR) agent observes a large file write to local storage. Neither surface
has instrumented the structural integrity of the artifact that was just
admitted into the environment.

Three weight file formats account for the substantial majority of model
artifacts in contemporary enterprise inference deployments. PyTorch
\texttt{.pt} and \texttt{.pth} checkpoints use Python's \texttt{pickle}
serialization protocol, which by design permits arbitrary Python objects to
be deserialized, creating a direct remote code execution (RCE) vector at
the moment \texttt{torch.load()} is called~\cite{pickleball2025}. Safetensors
was introduced specifically to eliminate deserialization-time code execution,
and it succeeds in that narrow objective; however, the format's JSON header,
which is parsed before any tensor data is loaded into memory, presents its
own attack surface through parser exploitation, metadata field abuse, and
the silent embedding of neural backdoor weights~\cite{safetensors2025empirical}.
GGUF, the dominant format for quantized inference with llama.cpp and Ollama,
includes a flexible key-value metadata section with no schema enforcement, and
multiple heap-buffer-overflow vulnerabilities in its reference parsing
implementation were disclosed by Cisco Talos in early 2024~\cite{gguf_quantization_attack2024}.

The practical attack pathway is straightforward. A threat actor forks a
popular model repository on Hugging Face, introduces a superficially plausible
commit, and publishes modified weight files. Engineers searching for a
specific model version frequently pull from forks without verifying
cryptographic provenance. The artifact enters the MLOps pipeline, clears
code-focused CI checks, and reaches a production inference endpoint. Hugging
Face does not cryptographically sign model artifacts by default. Hash values
published in model cards are informational metadata that any repository
owner, including a malicious one, can update after the fact.

This paper makes the following contributions:

\begin{itemize}
  \item A format-level threat model that characterizes the attack surface
    of PyTorch checkpoints, Safetensors archives, and GGUF binaries, grounded
    in documented incidents and peer-reviewed vulnerability research.
  \item \texttt{verify\_weights.py}, a Python binary inspector that performs
    Safetensors header anomaly detection and GGUF magic-byte and metadata
    key auditing without loading tensor data into memory.
  \item A five-stage GitHub Actions provenance gate
    (\texttt{model\_scan.yml}) that integrates registry hash verification,
    Sigstore/Cosign signature checking, \texttt{picklescan} opcode analysis,
    and binary inspection into a blocking CI/CD workflow.
  \item An empirical evaluation against eight synthetic adversarial test
    cases, demonstrating zero false negatives and zero false positives under
    controlled conditions.
  \item An open-source release of the complete toolkit at
    \url{https://github.com/sunilgentyala/model-provenance-guard}.
\end{itemize}

The remainder of the paper proceeds as follows. Section~\ref{sec:background}
surveys related work. Section~\ref{sec:threat} formalizes the threat model.
Section~\ref{sec:design} describes the system architecture. Section~\ref{sec:impl}
details the implementation. Section~\ref{sec:eval} presents the evaluation.
Section~\ref{sec:discussion} discusses limitations and future directions.
Section~\ref{sec:conclusion} concludes.

%% -----------------------------------------------------------------------
\section{Background and Related Work}
\label{sec:background}
%% -----------------------------------------------------------------------

\subsection{Machine Learning Supply Chain Attacks}

The concept of a software supply chain attack is well established in the
security literature, but its application to machine learning weight artifacts
is comparatively recent. Jiang et al.~\cite{modelsarecodes2024} conducted
a systematic measurement study of malicious code poisoning attacks on
pre-trained model hubs and found that attackers exploit the implicit trust
engineers place in popular repository namespaces. Separately, documented
incidents on Hugging Face have confirmed that malicious model artifacts can
execute reverse shell payloads at load time, yielding full control of
GPU-enabled inference infrastructure. The Communications of the ACM survey
by~\citet{maliciousai2025} argues that the combination of high download
volumes and weak provenance controls makes model hubs uniquely attractive
targets relative to traditional software package registries.

The zero-trust perspective on AI data pipelines, articulated by
\citet{mudusu2026zerotrust}, establishes that the never-trust-always-verify
principle must be applied not merely to network perimeters but directly to
data and model artifacts throughout the pipeline. That work demonstrates
that cryptographic authentication and tamper-evident lineage tracking,
implemented with ECDSA-P384 tokens and continuous policy enforcement, can
achieve complete detection of data tampering and component impersonation in
production AI systems. Our framework extends those principles specifically
to weight file artifacts at the CI/CD intake boundary.

\subsection{Pickle Deserialization in PyTorch Checkpoints}

The vulnerability class most thoroughly studied in the context of ML weight
files is pickle deserialization. Python's \texttt{pickle} module supports a
\texttt{\_\_reduce\_\_} protocol that enables arbitrary Python callables to
execute during deserialization. Any \texttt{torch.load()} call on a
maliciously constructed checkpoint can trigger operating system command
execution without any user interaction~\cite{pickleball2025}. Koishybayev
et al.~\cite{pickleball2025} present PickleBall, a secure deserialization
system that intercepts pickle opcodes before object reconstruction and
enforces a safe-by-default allowlist, demonstrating that the vulnerability
is addressable but requires a purpose-built defense rather than a general
virus scanner. JFrog Security Research additionally disclosed three zero-day
vulnerabilities in the widely deployed \texttt{picklescan} tool itself,
underscoring that even dedicated opcode scanners carry their own risk
surface.

\subsection{Safetensors Format and Header Security}

The Safetensors format was developed by Hugging Face as a direct response
to the RCE risks inherent in pickle-based serialization~\cite{safetensors2025empirical}.
By replacing pickle with a JSON-encoded header followed by raw tensor byte
data, Safetensors eliminates arbitrary code execution at load time. The
Trail of Bits security assessment of the Safetensors library~(2023)
confirmed the absence of deserialization vulnerabilities in the reference
implementation. However, the header parsing step introduces its own attack
surface: oversized JSON headers, malformed Unicode sequences, invalid tensor
dtype values, and data offsets that exceed file boundaries all represent
anomalies that warrant blocking behavior. CryptoTensors~\cite{cryptotensors2024}
proposes extending this format with tensor-level encryption and embedded
access control policies, indicating an active research community around
format security beyond the elimination of deserialization.

\subsection{GGUF Format Vulnerabilities}

GGUF, introduced by the \texttt{ggml} project in August 2023 and
subsequently adopted as the standard distribution format for quantized
models in the Ollama and llama.cpp ecosystems, presents a distinct set of
security concerns. Cisco Talos disclosed three heap-buffer-overflow
vulnerabilities in llama.cpp's GGUF parser in early 2024, all arising from
insufficient bounds checking in the metadata key-value parsing code. These
vulnerabilities confirm that GGUF files consumed by native C/C++
implementations can trigger memory corruption in unpatched consumers.
Beyond implementation flaws, researchers have demonstrated that quantization
error in GGUF models provides sufficient numerical flexibility to construct
adversarial models whose malicious behavior is invisible in full-precision
analysis but activates under quantized inference~\cite{gguf_quantization_attack2024}.

\subsection{Neural Backdoor Attacks}

The neural backdoor threat, formalized by \citeauthor{gu2019badnets} in the
seminal BadNets paper~\cite{gu2019badnets}, involves poisoning model weights
during training such that the model produces correct outputs on clean inputs
but triggers a targeted misclassification on inputs containing a specific
pattern. This attack class is particularly concerning because it is
undetectable by any binary-level inspection technique: the weight file loads
cleanly, passes format validation, and produces correct behavior under
standard test inputs. Detection requires behavioral analysis at inference
time, specifically monitoring for distributional anomalies in output
probabilities conditioned on statistical baselines. Binary inspection,
including the approach presented in this paper, explicitly does not address
neural backdoors; we characterize this as a residual risk and a direction
for complementary defense.

\subsection{Software Artifact Signing}

Sigstore~\cite{newman2022sigstore} provides a keyless signing infrastructure
that ties artifact signatures to ephemeral OIDC identity tokens, dramatically
lowering the adoption barrier relative to traditional GPG-based signing. The
Cosign tool, part of the Sigstore ecosystem, supports signing and verification
of arbitrary binary objects beyond container images, making it directly
applicable to model weight files. As of 2025, Sigstore underpins artifact
signing for npm, PyPI, and Kubernetes, and the OpenSSF AI/ML Working Group
has adapted these mechanisms specifically for model signing in a dedicated
specification~\cite{openssf_modelsigning2025}.

\subsection{ML Provenance and Governance}

Schlegel and Sattler~\cite{schlegel2024provenance} present a PROV-compliant
provenance model for end-to-end ML pipelines that captures artifact lineage
through MLflow and Git activities, enabling tamper-evident audit trails
across training and deployment stages. \citeauthor{gentyala2026aisbom}~\cite{gentyala2026aisbom}
frames the AI supply chain problem in terms of software bills of materials
(SBOMs) extended to AI components, demonstrating that SBOM standards such
as CycloneDX and SPDX can be adapted to encode model weight metadata,
training data references, and cryptographic attestations. The NIST AI Risk
Management Framework~\cite{nist_ai_rmf2023} and MITRE ATLAS~\cite{mitre_atlas2023}
provide complementary governance and threat classification scaffolding. The
OWASP Top 10 for LLM Applications explicitly names supply chain compromise
as a critical risk category~\cite{owasp_llm_supply2025}.

%% -----------------------------------------------------------------------
\section{Threat Model}
\label{sec:threat}
%% -----------------------------------------------------------------------

\subsection{Adversary Model}

We consider a threat actor who has obtained write access to a model
repository on a public model hub through one of three mechanisms: account
compromise of a legitimate model publisher, submission of a malicious pull
request that is merged by an inattentive or coerced maintainer, or creation
of a confusingly named fork designed to capture traffic from engineers who
do not verify repository provenance. The adversary's goal is to deliver a
weaponized weight artifact to at least one GPU-enabled inference endpoint
inside a target enterprise environment, either to execute arbitrary code at
model load time (for PyTorch checkpoints), to exploit parser vulnerabilities
in native runtimes (for GGUF), or to embed dormant backdoor weights that
activate on a trigger the adversary controls (applicable to all formats).

We assume the adversary has read-only access to the public model hub API
and can observe which model versions are popular and frequently downloaded.
We do not assume the adversary has network access to the target organization's
internal infrastructure. We assume the target organization's CI/CD pipeline
downloads model artifacts automatically from external sources as part of its
standard model update workflow, which is the common MLOps pattern.

\subsection{PyTorch Checkpoint Attack Surface}

PyTorch checkpoints use Python's \texttt{pickle} serialization protocol.
The \texttt{pickle} module supports a \texttt{\_\_reduce\_\_} method that,
when present on a serialized object, is called by the deserializer to
reconstruct the object. An attacker can craft a checkpoint that includes an
object whose \texttt{\_\_reduce\_\_} method returns a callable paired with
arguments, effectively embedding an arbitrary function call that executes
the moment \texttt{torch.load()} is invoked. Common payloads include reverse
shells, credential harvesters targeting cloud provider metadata endpoints,
and persistence mechanisms that install themselves into the Python runtime.
No user interaction beyond the standard model loading call is required.

\subsection{Safetensors Attack Surface}

While the Safetensors format eliminates deserialization-time code execution,
it introduces three distinct attack surfaces. First, the JSON header
processed before tensor data is loaded is vulnerable to parser exploits:
oversized headers exceeding implementation-specific memory limits, embedded
null bytes that confuse string parsers, and deeply nested JSON structures
that trigger stack exhaustion in recursive parsers all represent viable
denial-of-service vectors. Second, the \texttt{\_\_metadata\_\_} dictionary
accepts arbitrary string key-value pairs with no schema enforcement;
maliciously crafted values in this field can exploit applications that log,
display, or forward metadata without sanitization. Third, and most
consequentially for high-value enterprise targets, Safetensors files can
carry weights poisoned at training time~\cite{gu2019badnets}. The file loads
cleanly under all format-level checks. The backdoor activates only on
attacker-chosen inputs.

\subsection{GGUF Attack Surface}

GGUF's binary metadata section has no schema enforcement beyond type tagging.
A threat actor can embed arbitrary data under custom metadata keys; in native
C and C++ consumers, keys or values whose lengths exceed what the consuming
code anticipates can trigger heap-based buffer overflows. The version field
at bytes 4--7 controls how the remainder of the file is parsed; supplying an
unsupported version value can trigger undefined behavior in parsers that do
not include explicit version gating. The tensor count field at bytes 8--15
is trusted by consuming libraries; a declared count that differs from the
number of tensors actually present in the file causes the parser to
over-read memory.

\subsection{Practical Attack Path}

The end-to-end attack path requires no novel infrastructure and is
reproducible by a moderately skilled threat actor. The actor creates a
Hugging Face account with a name that differs from a popular model
publisher's account by a single character transposition. They publish a
modified checkpoint whose malicious payload is embedded in a tensor that is
never accessed during standard evaluation. Engineers performing searches for
the original model encounter the forked repository, note comparable download
statistics and a convincing model card, and pull the artifact. The artifact
clears all code-focused CI checks because no ML model inspection is
configured. It is registered in the internal model registry, promoted to
staging, and eventually deployed to a production inference endpoint, at
which point the payload executes.

%% -----------------------------------------------------------------------
\section{System Architecture}
\label{sec:design}
%% -----------------------------------------------------------------------

\subsection{Design Principles}

The model-provenance-guard framework is grounded in four design principles
derived from established supply chain security practice and the zero-trust
architecture posture described by~\citet{mudusu2026zerotrust}.

\textbf{Shift-left enforcement.} Provenance checks must occur at the model
intake boundary, before any artifact is registered in an internal model
registry or promoted beyond the staging environment. Checks that operate
after deployment provide no prevention value.

\textbf{Defense in depth.} No single control is sufficient. Hash
verification catches substitution attacks but not artifacts that were never
in the registry. Signature verification catches unsigned artifacts but not
ones signed by a compromised key. Binary inspection catches format-level
anomalies but not neural backdoors. The pipeline combines all four controls
because each addresses a distinct portion of the threat surface.

\textbf{Hard failures, not warnings.} Any control failure must block pipeline
progression. A provenance gate that emits warnings while allowing promotion
to continue provides no practical security benefit and creates compliance
theater.

\textbf{Registry integrity bootstrapping.} The trusted hash registry must
itself be protected against tampering. The pipeline validates the registry's
own SHA-256 hash against a value stored as a sealed repository secret before
accepting any registry lookup result.

\subsection{Pipeline Architecture}

The provenance gate is implemented as a GitHub Actions workflow with five
sequentially dependent jobs, as illustrated in Figure~\ref{fig:pipeline}.

\begin{figure}[h]
\centering
\fbox{\parbox{0.85\columnwidth}{%
\textbf{Job 1: Registry Integrity Check}\\
Computes SHA-256 of trusted\_models.json and\\
compares against a sealed repository secret.\\[4pt]
\textbf{Job 2: Hash Verification}\\
Downloads artifact, computes SHA-256, and checks\\
against the trusted registry (not the model card).\\[4pt]
\textbf{Job 3: Cosign Signature Verification}\\
Verifies .bundle or .sig file against an organization\\
public key stored as a repository secret.\\[4pt]
\textbf{Job 4: Binary Format Inspection}\\
Routes to picklescan (PyTorch) or verify\_weights.py\\
(Safetensors / GGUF) based on file extension.\\[4pt]
\textbf{Job 5: Provenance Gate}\\
Evaluates all upstream results. Fails the pipeline\\
if any required job did not succeed.
}}
\caption{Five-stage provenance gate pipeline.}
\label{fig:pipeline}
\end{figure}

The workflow triggers on three event types: pull requests that modify any
path containing model artifact files (\texttt{*.safetensors},
\texttt{*.gguf}, \texttt{*.pt}, \texttt{*.pth}), a nightly schedule that
scans the organization's staging model cache for drift, and
\texttt{workflow\_dispatch} for on-demand verification of specific artifacts.

\subsection{Trusted Registry Design}

The trusted model registry is a JSON file maintained in the same repository
as the workflow, protected by branch protection rules requiring at least
one approved reviewer on any pull request that modifies it. Each entry
records the model name, the artifact filename, the SHA-256 hash computed by
a human reviewer at initial review time, the source URL and source repository
commit hash, the reviewing engineer's identity, the review timestamp, and an
optional Cosign bundle path.

A critical design decision is that the trusted hash must be computed by a
human reviewer from the artifact at initial review time and must not be
sourced from the model card or model hub metadata. An adversary with control
over a model repository can update the published hash at any time. Only
hashes computed independently and stored in a separately controlled registry
can provide supply chain integrity guarantees.

%% -----------------------------------------------------------------------
\section{Implementation}
\label{sec:impl}
%% -----------------------------------------------------------------------

\subsection{Safetensors Header Inspector}

The \texttt{inspect\_safetensors()} function performs seven checks before
returning a pass or fail verdict, and in no case does it load tensor data
into memory.

\textbf{Check 1: Minimum file size.} The Safetensors specification requires
the first eight bytes to encode the header length as a little-endian
unsigned 64-bit integer. A file smaller than eight bytes cannot be a valid
Safetensors file and is rejected immediately.

\textbf{Check 2: Header size bounds.} The declared header size is validated
against a configurable ceiling (100~MB by default) and a warning threshold
(10~MB). Headers exceeding the ceiling indicate a possible header injection
attack or a severely malformed file. Headers between the warning threshold
and ceiling are flagged for manual review.

\textbf{Check 3: UTF-8 validity.} Header bytes must decode as valid UTF-8
before JSON parsing is attempted. Files that fail UTF-8 decoding are rejected
before any JSON parser is invoked, preventing a class of attacks that exploit
parser behavior on malformed Unicode sequences.

\textbf{Check 4: JSON validity.} The decoded header string must parse as
valid JSON. A malformed JSON header is treated as a hard anomaly rather than
a soft warning because legitimate Safetensors files produced by any major
framework always contain well-formed JSON headers.

\textbf{Checks 5--6: Tensor entry validation.} Each tensor entry in the
header must contain a \texttt{dtype} field drawn from the set of valid
Safetensors data types (\texttt{F64}, \texttt{F32}, \texttt{F16},
\texttt{BF16}, \texttt{I64}, \texttt{I32}, \texttt{I16}, \texttt{I8},
\texttt{U8}, \texttt{BOOL}), a well-formed integer-list \texttt{shape} field
with non-negative dimensions, and a two-element integer-list
\texttt{data\_offsets} field with a non-negative start and an end no less
than the start. Any entry failing these checks is classified as anomalous.

\textbf{Check 7: Data offset bounds.} The maximum data offset declared
across all tensor entries is added to the header region's starting position
and compared against the file size. A declared data region that extends
beyond the end of the file is treated as a hard failure. This check catches
truncated files and files constructed to declare phantom data regions as a
precursor to parser exploitation.

\subsection{GGUF Inspector}

The \texttt{inspect\_gguf()} function validates the binary header and
audits metadata key names without parsing metadata values, which would
require a full GGUF parser and introduce its own attack surface.

The function first checks that the file begins with the four magic bytes
\texttt{GGUF} (ASCII). It then reads the version field, flagging versions
outside the set \{2, 3\} as anomalous while permitting inspection to
continue (unless \texttt{GGUF\_STRICT=1} is set in the environment). It
reads the declared tensor count and key-value count, flagging the latter
if it exceeds 10,000, which far exceeds any count observed in production
models from major providers.

For each metadata key, the function reads the key length as a
little-endian 64-bit integer, rejects keys longer than 1,024 bytes, reads
the key bytes, validates UTF-8 encoding, and checks whether the key name
begins with one of 28 known architecture-specific prefixes defined in the
GGUF v3 specification (\texttt{general.}, \texttt{llama.},
\texttt{mistral.}, \texttt{tokenizer.}, etc.). Keys that do not match any
expected prefix are flagged as anomalous. Because the GGUF format
interleaves key names with variable-length, variable-type values, fully
sequential key parsing without a complete value-skip implementation is not
safe; the current inspector reads the first key and classifies it, with a
documented path to enable full sequential parsing using a complete
GGUF value-skip implementation.

\subsection{Cross-Platform Correctness}

The initial implementation hardcoded \texttt{/tmp} as the report output
directory and used \texttt{datetime.utcnow()}, which was deprecated in
Python~3.12. We corrected the report path to use \texttt{tempfile.gettempdir()},
which resolves correctly on Windows, Linux, and macOS, and replaced the
deprecated call with \texttt{datetime.now(datetime.timezone.utc)}.

\subsection{Hash Verification and Registry Lookup}

The \texttt{verify\_against\_registry()} function computes the SHA-256
digest of the artifact in 1~MB chunks to accommodate files that routinely
exceed 4~GB, then queries the trusted registry for a matching entry by
filename. If no entry is found, the function returns a permissive result
with a warning, leaving the hard failure to the GitHub Actions pipeline
step, which is configured to treat an unregistered artifact as a blocking
error. If an entry is found but the hash does not match, the function
returns a hard failure immediately.

%% -----------------------------------------------------------------------
\section{Evaluation}
\label{sec:eval}
%% -----------------------------------------------------------------------

\subsection{Test Methodology}

We constructed eight synthetic test cases covering the Safetensors and GGUF
formats. Each test case programmatically constructs a binary artifact that
embodies a specific structural condition---either a valid format or a
specific anomaly---and invokes the inspector with the \texttt{--skip-hash}
flag to isolate the binary inspection controls from the registry hash check.
The expected exit code (0 for pass, 1 for failure) is compared against the
actual exit code to determine test outcome.

PyTorch checkpoint testing is delegated to \texttt{picklescan}, a
well-validated tool with its own test suite; we do not redundantly test
functionality that picklescan already covers. Neural backdoor detection is
explicitly outside the scope of binary inspection and is not tested.

\subsection{Results}

Table~\ref{tab:results} summarizes the outcomes across all eight test cases.

\begin{table}[h]
\caption{Binary inspector evaluation results.}
\label{tab:results}
\begin{tabular}{p{0.52\columnwidth}cc}
\toprule
\textbf{Test Case} & \textbf{Expected} & \textbf{Result} \\
\midrule
Safetensors: valid file (F32, correct offsets)  & PASS & \textbf{PASS} \\
Safetensors: invalid dtype (\texttt{EVIL\_TYPE}) & FAIL & \textbf{FAIL} \\
Safetensors: data offset past EOF               & FAIL & \textbf{FAIL} \\
Safetensors: non-JSON header bytes              & FAIL & \textbf{FAIL} \\
GGUF: valid file, zero KV entries               & PASS & \textbf{PASS} \\
GGUF: valid \texttt{general.architecture} key   & PASS & \textbf{PASS} \\
GGUF: wrong magic bytes (\texttt{EVIL})         & FAIL & \textbf{FAIL} \\
GGUF: unsupported version (v99), non-strict mode & PASS & \textbf{PASS} \\
\bottomrule
\end{tabular}
\end{table}

All eight test cases produced the expected outcome, yielding a false
negative rate of 0.0\% and a false positive rate of 0.0\% under the tested
conditions. The unsupported GGUF version case correctly produces a warning
annotation in the inspection report while allowing the artifact through,
consistent with the design intent that version anomalies should surface for
human review without blocking benign artifacts from non-standard but
structurally sound files.

\subsection{Performance Characteristics}

Binary inspection latency is bounded by file read bandwidth rather than
computation. The Safetensors inspector reads only the header region, so
inspection time for a 7~GB model file is dominated by the time to seek to
and read the header, typically under two seconds on standard SSD storage.
The GGUF inspector reads the fixed-length header and the first metadata key,
a region that is always well under 1~KB, making its latency effectively
constant regardless of file size. Hash computation scales linearly with file
size; at 1~MB chunk reads on modern storage, a 7~GB file completes in
approximately 14~seconds.

%% -----------------------------------------------------------------------
\section{Discussion}
\label{sec:discussion}
%% -----------------------------------------------------------------------

\subsection{Limitations}

The most consequential limitation of the framework is its inability to
detect neural backdoors embedded in weight values. A Safetensors or GGUF
file whose weights have been poisoned at training time passes all structural
checks cleanly, loads without incident, and behaves correctly on all clean
inputs. The malicious behavior surfaces only on inputs containing the
attacker's trigger pattern. Addressing this requires behavioral monitoring
at inference time: continuous analysis of output distributions, activation
statistics, and response latency to identify anomalies that deviate from
a baseline established on known-clean model weights~\cite{mitre_atlas2023}.
This is a complementary defense that operates at the application layer, not
the binary layer, and it falls outside the scope of the present framework.

The GGUF inspector's metadata audit is currently limited to key name
classification. Full value inspection, which would enable detection of
oversized values and values whose types are inconsistent with the
architecture specification, requires a complete GGUF value-skip implementation
that correctly advances the file pointer across all GGUF value types. We
document this limitation explicitly and recommend that teams requiring full
metadata value inspection use \texttt{gguf-py} or llama.cpp's
\texttt{gguf-dump} utility as a complement.

\subsection{Scope and Integration}

The framework is designed to sit at the model intake boundary: between the
artifact download step and the step that registers the artifact in an
internal model registry. It does not address security controls at earlier
stages (model training integrity, dataset provenance) or at later stages
(inference-time behavioral monitoring). Organizations should treat this
framework as one layer in a defense-in-depth strategy aligned with the NIST
AI RMF~\cite{nist_ai_rmf2023} and the OWASP LLM Top~10~\cite{owasp_llm_supply2025}.

\subsection{Future Work}

Four extensions represent near-term priorities. First, extending the binary
inspector to cover ONNX protobuf models and TensorFlow SavedModel artifacts,
which introduce their own format-specific attack surfaces not addressed by
the current implementation. Second, integrating AI Software Bill of
Materials (AI-SBOM) generation into the provenance gate output, enabling
downstream systems to consume cryptographically attested lineage records in
standard SBOM format~\cite{gentyala2026aisbom,frontiers2026aibom}. Third,
implementing complete GGUF metadata value parsing to enable schema-level
validation beyond key name auditing. Fourth, investigating integration with
inference-time behavioral monitoring to close the neural backdoor detection
gap, for example by coupling the provenance gate output with a downstream
monitoring hook that enforces output distribution baselines.

%% -----------------------------------------------------------------------
\section{Conclusion}
\label{sec:conclusion}
%% -----------------------------------------------------------------------

The machine learning supply chain is not merely a software supply chain. It
is a binary supply chain in which the most sensitive components---weight
files that encode model behavior---are distributed through infrastructure
that provides no cryptographic integrity guarantees by default and consumed
by tooling that offers no structural inspection prior to loading. This gap
is neither obscure nor theoretical; it has been exploited in documented
incidents and characterized in peer-reviewed vulnerability research across
all three dominant weight file formats.

We have presented model-provenance-guard, a framework that addresses this
gap with a combination of binary-level inspection and cryptographic
provenance enforcement, deployed at the CI/CD boundary where prevention
remains possible. The framework requires no architectural changes to
existing MLOps pipelines; it inserts as a gate between artifact download
and registry registration. Against eight synthetic adversarial test cases
spanning the Safetensors and GGUF attack surfaces, the inspector returned
zero false negatives and zero false positives.

The underlying principle is not novel: build the gate, not more alerts.
Monitoring a threat class that cannot be observed through existing
instrumentation generates no signal. The response to weaponized weights is
structural, and it must be implemented before the next artifact a team
stages for production introduces a threat that no existing tool was designed
to see.

%% -----------------------------------------------------------------------
\begin{acks}
The authors thank the security research community for disclosing the
llama.cpp heap-buffer-overflow vulnerabilities (Cisco Talos, 2024) and the
Trail of Bits team for the Safetensors library security assessment (2023),
both of which directly informed the threat model developed in this work.
The model-provenance-guard toolkit is released under the MIT License and
is available at \url{https://github.com/sunilgentyala/model-provenance-guard}.
\end{acks}

%% -----------------------------------------------------------------------
\balance
\bibliographystyle{ACM-Reference-Format}
\bibliography{references}

\end{document}
